\begin{lemma}[Failure and legal moves in the copied triangle]
Let $h$ be bad for $\mathsf{PowerRel}(r)$.  For every increasing
enumeration $x$, row $n$, and depth $k$, the copied pair is failed and the
next left value is a legal move:
\[
\begin{split}
&\neg\mathsf{PowerRel}(r)(L_{n,k},L_{n+1,k}),\\
&\mathsf{PowerMove}(L_{n,k},L_{n,k+1}),
\end{split}
\]
where $L_{n,k}=\mathsf{PowerCopiedLeft}_r(h,x,n,k)$.
\end{lemma}
\begin{proof}
Fix an arbitrary increasing enumeration $x$.  Using the notation of
\noderef{PowerCopiedLeft}, write
\[
 S_i=\mathsf{PowerShiftIter}(i,x)
 \quad\text{and}\quad
 L_{i,j}=\mathsf{PowerCopiedLeft}_r(h,x,i,j).
\]
We prove simultaneously for every row $i$ and every depth $j$ that
\[
 F_{i,j}:=\neg\mathsf{PowerRel}(r)(L_{i,j},L_{i+1,j})
 \quad\text{and}\quad
 M_{i,j}:=\mathsf{PowerMove}(L_{i,j},L_{i,j+1}).
\]

We use induction on the depth.  First consider depth zero.  By the definition of badness in
\noderef{BadMultiSequence}, applied to the increasing enumeration $S_i$,
we have
\[
 \neg\mathsf{PowerRel}(r)(h(S_i),h(\mathsf{ShiftMap}(S_i))).
 \tag{1}
\]
The successor equation in \noderef{PowerShiftIter} identifies
$\mathsf{ShiftMap}(S_i)$ with $S_{i+1}$.  Thus, for every $i$,
\[
 R_i:=\neg\mathsf{PowerRel}(r)(h(S_i),h(S_{i+1})).
 \tag{2}
\]
In particular, (2) gives the raw failures $R_n$ and $R_{n+1}$ for the two
adjacent rows under consideration.

The depth-zero clause of \noderef{PowerCopiedLeft} gives, for every $i$,
\[
 L_{i,0}=\mathsf{PowerIResponse}_r(h(S_i),h(S_{i+1})).
 \tag{3}
\]
Applying \noderef{PowerIResponseLegal} to $R_n$ and using (3) gives the
legal move
\[
 \mathsf{PowerMove}(h(S_n),L_{n,0}).
 \tag{4}
\]
Applying the same lemma to $R_{n+1}$ and using (3) at row $n+1$ gives
\[
 \mathsf{PowerMove}(h(S_{n+1}),L_{n+1,0}).
 \tag{5}
\]
Now apply \noderef{PowerIResponseFailure} to the raw failure $R_n$ and the
right-hand legal move (5).  Its conclusion is
\[
 \neg\mathsf{PowerRel}(r)\bigl(
 \mathsf{PowerIResponse}_r(h(S_n),h(S_{n+1})),L_{n+1,0}\bigr).
 \tag{6}
\]
Using (3) once more, (6) is exactly $F_{n,0}$.

Finally apply \noderef{PowerIResponseLegal} to the failed pair in (6), after
rewriting its first coordinate by (3).  The successor-depth clause of
\noderef{PowerCopiedLeft} says
\[
 L_{n,1}=\mathsf{PowerIResponse}_r(L_{n,0},L_{n+1,0}).
 \tag{7}
\]
Consequently the resulting legal response is
$\mathsf{PowerMove}(L_{n,0},L_{n,1})$, which is $M_{n,0}$.

For the induction step, assume that $F_{i,k}$ and $M_{i,k}$ hold for every
row $i$.  In particular, $F_{n,k}$ holds, and $M_{n+1,k}$ gives the legal
right-row move
\[
 \mathsf{PowerMove}(L_{n+1,k},L_{n+1,k+1}).
 \tag{8}
\]
The successor-depth clause of \noderef{PowerCopiedLeft}, applied once to
row $n$ and once to row $n+1$, gives the two equations
\[
 \begin{split}
 L_{n,k+1}&=\mathsf{PowerIResponse}_r(L_{n,k},L_{n+1,k}),\\
 L_{n+1,k+1}&=\mathsf{PowerIResponse}_r(L_{n+1,k},L_{n+2,k}).
 \end{split}
 \tag{9}
\]
Apply \noderef{PowerIResponseFailure} to $F_{n,k}$ and the legal move (8).
It yields
\[
 \neg\mathsf{PowerRel}(r)\bigl(
 \mathsf{PowerIResponse}_r(L_{n,k},L_{n+1,k}),L_{n+1,k+1}\bigr).
 \tag{10}
\]
The first equation in (9) rewrites (10) as $F_{n,k+1}$.

Apply \noderef{PowerIResponseLegal} to the failed pair (10).  The
successor-depth clause of \noderef{PowerCopiedLeft}, now at depth $k+1$,
also gives
\[
 L_{n,k+2}=\mathsf{PowerIResponse}_r(L_{n,k+1},L_{n+1,k+1}).
 \tag{11}
\]
Using the first equation in (9) and then (11), the legal response supplied
by \noderef{PowerIResponseLegal} is precisely
\[
 \mathsf{PowerMove}(L_{n,k+1},L_{n,k+2}),
\]
which is $M_{n,k+1}$.  Thus both assertions hold at depth $k+1$ for row
$n$, and since $n$ was arbitrary the induction is simultaneous over all
rows.  Therefore the two claims in the lemma hold for every $n$ and $k$.
Since $x$ was arbitrary, this proves the lemma for every increasing
enumeration.
\end{proof}
